% answers.tex —— 2014 年·广东卷·理科 · 参考答案与详细解析片段
% 由 tools/generate_projects.py 依来源总账自 2014/2014_数学真题_广东理_试题.md + 2014/2014_数学真题_广东理_解析正文.md + 2014/2014_数学真题_广东理_答案要点.md 生成，勿手改（改源 MD 后重跑）


\begin{answer}
\gksolhead{1}{}

\textbf{题目}：

已知集合 $M=\{-1,0,1\}$，$N=\{0,1,2\}$，则 $M\cup N=$\mbox{（　　）}

（A）$\{0,1\}$；（B）$\{-1,0,2\}$；（C）$\{-1,0,1,2\}$；（D）$\{-1,0,1\}$。

\Needspace{4\baselineskip}
\textbf{解析}

\Needspace{5\baselineskip}
$M=\{-1,0,1\}$，$N=\{0,1,2\}$，并集就是两集合的全部元素，去重后为
\nopagebreak[4]
\[M\cup N=\{-1,0,1,2\}.\]

选 C。

\gkanswer{答案：C。$M\cup N=\{-1,0,1\}\cup\{0,1,2\}=\{-1,0,1,2\}$。}

\end{answer}
\begin{answer}
\gksolhead{2}{}

\textbf{题目}：

已知复数 $z$ 满足 $(3+4\mathrm i)z=25$，则 $z=$\mbox{（　　）}

（A）$-3+4\mathrm i$；（B）$-3-4\mathrm i$；（C）$3+4\mathrm i$；（D）$3-4\mathrm i$。

\Needspace{4\baselineskip}
\textbf{解析}

\Needspace{5\baselineskip}
两边同乘 $3-4\mathrm i$，利用 $(3+4\mathrm i)(3-4\mathrm i)=9+16=25$，得
\nopagebreak[4]
\[z=\frac{25}{3+4\mathrm i}=\frac{25(3-4\mathrm i)}{25}=3-4\mathrm i.\]

选 D。

\gkanswer{答案：D。$z=\dfrac{25}{3+4\mathrm i}=\dfrac{25(3-4\mathrm i)}{25}=3-4\mathrm i$。}

\end{answer}
\begin{answer}
\gksolhead{3}{}

\textbf{题目}：

若变量 $x$，$y$ 满足约束条件 $\begin{cases}y\leqslant x,\\x+y\leqslant 1,\\y\geqslant -1,\end{cases}$ 且 $z=2x+y$ 的最大值和最小值分别为 $m$ 和 $n$，则 $m-n=$\mbox{（　　）}

（A）$5$；（B）$6$；（C）$7$；（D）$8$。

\Needspace{4\baselineskip}
\textbf{解析}

\Needspace{5\baselineskip}
约束条件 $y\leqslant x$，$x+y\leqslant 1$，$y\geqslant -1$ 围成三角形，三边两两相交得三个顶点：
\nopagebreak[4]
\[y=x\ \text{\text{与}}\ y=-1\ \Rightarrow\ (-1,-1),\]
\nopagebreak[4]
\[y=x\ \text{\text{与}}\ x+y=1\ \Rightarrow\ \left(\frac12,\frac12\right),\]
\nopagebreak[4]
\[x+y=1\ \text{\text{与}}\ y=-1\ \Rightarrow\ (2,-1).\]

\Needspace{5\baselineskip}
把三个顶点代入 $z=2x+y$：
\nopagebreak[4]
\[\begin{aligned}
(-1,-1)&:\ z=2\times(-1)+(-1)=-3,\\
\left(\frac12,\frac12\right)&:\ z=2\times\frac12+\frac12=\frac32,\\
(2,-1)&:\ z=2\times2+(-1)=3.
\end{aligned}\]
\Needspace{5\baselineskip}
线性目标函数在三角形可行域上的最值必在顶点处取得，故
\nopagebreak[4]
\[m=3,\qquad n=-3,\qquad m-n=3-(-3)=6.\]

选 B。

\gkanswer{答案：B。可行域顶点 $(-1,-1)$，$\left(\dfrac12,\dfrac12\right)$，$(2,-1)$ 处目标值依次为 $-3$，$\dfrac32$，$3$，故 $m=3$，$n=-3$，$m-n=6$。}

\end{answer}
\begin{answer}
\gksolhead{4}{}

\textbf{题目}：

若实数 $k$ 满足 $0<k<9$，则曲线 $\dfrac{x^2}{25}-\dfrac{y^2}{9-k}=1$ 与曲线 $\dfrac{x^2}{25-k}-\dfrac{y^2}{9}=1$ 的\mbox{（　　）}

（A）焦距相等；（B）实半轴长相等；（C）虚半轴长相等；（D）离心率相等。

\Needspace{4\baselineskip}
\textbf{解析}

\Needspace{5\baselineskip}
第一条曲线中 $a_1^2=25$，$b_1^2=9-k$，故
\nopagebreak[4]
\[c_1^2=a_1^2+b_1^2=25+(9-k)=34-k.\]
\Needspace{5\baselineskip}
第二条曲线中 $a_2^2=25-k$，$b_2^2=9$（因 $0<k<9$，两根号内均为正），故
\nopagebreak[4]
\[c_2^2=a_2^2+b_2^2=(25-k)+9=34-k.\]
于是 $c_1^2=c_2^2$，两曲线焦距相等。

再排除其余三项：实半轴长分别为 $5$ 与 $\sqrt{25-k}$，因 $0<k<9$ 而不等；虚半轴长分别为 $\sqrt{9-k}$ 与 $3$，不等；离心率分别为 $\dfrac{\sqrt{34-k}}{5}$ 与 $\dfrac{\sqrt{34-k}}{\sqrt{25-k}}$，分母不同，不等。

选 A。

\gkanswer{答案：A。两曲线的 $c^2$ 都为 $25+(9-k)=(25-k)+9=34-k$，故焦距相等。}

\end{answer}
\begin{answer}
\gksolhead{5}{}

\textbf{题目}：

已知向量 $\boldsymbol a=(1,0,-1)$，则下列向量中与 $\boldsymbol a$ 成 $60^\circ$ 夹角的是\mbox{（　　）}

（A）$(-1,1,0)$；（B）$(1,-1,0)$；（C）$(0,-1,1)$；（D）$(-1,0,1)$。

\Needspace{4\baselineskip}
\textbf{解析}

$|\boldsymbol a|=\sqrt{1^2+0^2+(-1)^2}=\sqrt2$。逐项计算夹角余弦：

（A）$\boldsymbol b=(-1,1,0)$，$\boldsymbol a\cdot\boldsymbol b=-1$，$|\boldsymbol b|=\sqrt2$，故 $\cos\langle\boldsymbol a,\boldsymbol b\rangle=-\dfrac12$，夹角 $120^\circ$；

（B）$\boldsymbol b=(1,-1,0)$，$\boldsymbol a\cdot\boldsymbol b=1$，$|\boldsymbol b|=\sqrt2$，故 $\cos\langle\boldsymbol a,\boldsymbol b\rangle=\dfrac12$，夹角 $60^\circ$；

（C）$\boldsymbol b=(0,-1,1)$，$\boldsymbol a\cdot\boldsymbol b=-1$，$|\boldsymbol b|=\sqrt2$，故 $\cos\langle\boldsymbol a,\boldsymbol b\rangle=-\dfrac12$，夹角 $120^\circ$；

（D）$\boldsymbol b=(-1,0,1)$，$\boldsymbol a\cdot\boldsymbol b=-2$，$|\boldsymbol b|=\sqrt2$，故 $\cos\langle\boldsymbol a,\boldsymbol b\rangle=-1$，夹角 $180^\circ$。

只有（B）成 $60^\circ$。

选 B。

\gkanswer{答案：B。$\boldsymbol a\cdot(1,-1,0)=1$，$|\boldsymbol a|=|(1,-1,0)|=\sqrt2$，夹角余弦为 $\dfrac12$，即 $60^\circ$。}

\end{answer}
\begin{answer}
\gksolhead{6}{}

\textbf{题目}：

已知某地区中小学生人数和近视情况如图 1 和图 2 所示，为了解该地区中小学生的近视形成原因，用分层抽样的方法抽取 $2\%$ 的学生进行调查，则样本容量和抽取的高中生近视人数分别为\mbox{（　　）}

（A）$200$，$20$；（B）$100$，$20$；（C）$200$，$10$；（D）$100$，$10$。
\nopagebreak[4]
\begin{center}% ============================================================
% f2014_guangdong_science_q06_charts.tex
% 2014 广东卷（理科）第 6 题：某地区中小学生人数饼图（图 1）与
% 各学段近视率柱状图（图 2）。
% 来源 PDF：普通高考/2014/2014广东理.pdf
% 原页：第 1 页（A3 横排）第 1 栏，印面 877，第 6 题题干下方
% 说明：图 1 中 小学生 3500 名、初中生 4500 名、高中生 2000 名，
% 三扇区圆心角依次为 126°、162°、72°；图 2 中 小学、初中、高中
% 的近视率依次为 10%、30%、50%。
% ============================================================
\begin{tikzpicture}[>=Stealth, line width=0.7pt, font=\small,
  x=1cm, y=1cm, line cap=round, line join=round]

  % ---------------- 图 1：人数饼图 ----------------
  \coordinate (O) at (0,0.9);
  \draw (O) circle (2);
  \draw (O) -- ++(90:2);
  \draw (O) -- ++(216:2);
  \draw (O) -- ++(18:2);

  \node[align=center, font=\footnotesize] at (-0.998,1.409)
    {小学生\\[1pt]$3500$ 名};
  \node[align=center, font=\footnotesize] at (0.509,-0.098)
    {初中生\\[1pt]$4500$ 名};
  \node[align=center, font=\footnotesize] at (0.658,1.790)
    {高中生\\[1pt]$2000$ 名};

  \node at (0,-1.5) {图 1};

  % ---------------- 图 2：近视率柱状图 ----------------
  \coordinate (B0) at (7,0);
  \draw[->] (B0) -- (11.9,0) node[below, inner sep=4pt] {年级};
  \draw[->] (B0) -- (7,4.25) node[above] {近视率/$\%$};

  % y 轴刻度与虚线
  \foreach \y/\lab/\xend in {0.9/10/8.3, 2.4/30/9.5, 3.6/50/10.7} {
    \draw (6.88,\y) -- (7.12,\y);
    \node[left, font=\footnotesize] at (6.82,\y) {$\lab$};
    \draw[dashed, line width=0.5pt] (7,\y) -- (\xend,\y);
  }
  \node[below left, font=\footnotesize, inner sep=3pt] at (7,0) {$O$};

  % 三根柱子
  \draw (7.6,0) rectangle (8.3,0.9);
  \draw (8.8,0) rectangle (9.5,2.4);
  \draw (10.0,0) rectangle (10.7,3.6);

  \node[below, font=\footnotesize] at (7.95,-0.06) {小学};
  \node[below, font=\footnotesize] at (9.15,-0.06) {初中};
  \node[below, font=\footnotesize] at (10.35,-0.06) {高中};

  \node at (9.15,-1.5) {图 2};
\end{tikzpicture}
\end{center}
\nopagebreak[4]
\Needspace{4\baselineskip}
\textbf{解析}

\Needspace{5\baselineskip}
由图 1，该地区中小学生总人数为
\nopagebreak[4]
\[3500+4500+2000=10000.\]
\Needspace{5\baselineskip}
按 $2\%$ 分层抽样，样本容量为
\nopagebreak[4]
\[10000\times 2\%=200.\]
\Needspace{5\baselineskip}
高中生共 $2000$ 人，抽中
\nopagebreak[4]
\[2000\times 2\%=40\]
\Needspace{5\baselineskip}
人；由图 2，高中生近视率为 $50\%$，故抽取的高中生中近视人数为
\nopagebreak[4]
\[40\times 50\%=20.\]

所以样本容量与抽取的高中生近视人数分别为 $200$，$20$。

选 A。

\gkanswer{答案：A。总人数 $10000$，样本容量 $10000\times2\%=200$；高中生抽 $2000\times2\%=40$ 人，近视率 $50\%$，故近视 $20$ 人。}

\end{answer}
\begin{answer}
\gksolhead{7}{}

\textbf{题目}：

若空间中四条两两不同的直线 $l_1$，$l_2$，$l_3$，$l_4$，满足 $l_1\perp l_2$，$l_2\perp l_3$，$l_3\perp l_4$，则下列结论一定正确的是\mbox{（　　）}

（A）$l_1\perp l_4$；（B）$l_1\parallel l_4$；（C）$l_1$，$l_4$ 既不垂直也不平行；（D）$l_1$，$l_4$ 的位置关系不确定。

\Needspace{4\baselineskip}
\textbf{解析}

题设只给出三次垂直，没有给出任何平行关系。取 $l_2$ 为 $z$ 轴、$l_3$ 为 $x$ 轴，则 $l_1\perp l_2$ 只要求 $l_1$ 的方向在 $xy$ 平面内，$l_4\perp l_3$ 只要求 $l_4$ 的方向在 $yz$ 平面内，二者方向可任意选取。

若取 $l_1$ 为 $y$ 轴、$l_4$ 为平行于 $y$ 轴的另一条直线，则 $l_1\parallel l_4$；若取 $l_1$ 为 $y$ 轴、$l_4$ 为平行于 $z$ 轴的直线，则 $l_1\perp l_4$；也可取成既不平行也不垂直的异面直线。三种位置关系都能出现，故 $l_1$，$l_4$ 的位置关系不确定。

选 D。

\gkanswer{答案：D。三次垂直不含任何平行条件，$l_1$ 与 $l_4$ 可平行、可垂直、也可异面，位置关系不确定。}

\end{answer}
\begin{answer}
\gksolhead{8}{}

\textbf{题目}：

设集合 $A=\{(x_1,x_2,x_3,x_4,x_5)\mid x_i\in\{-1,0,1\},\ i=1,2,3,4,5\}$，那么集合 $A$ 中满足条件“$1\leqslant|x_1|+|x_2|+|x_3|+|x_4|+|x_5|\leqslant 3$”的元素个数为\mbox{（　　）}

（A）$60$；（B）$90$；（C）$120$；（D）$130$。

\Needspace{4\baselineskip}
\textbf{解析}

$|x_i|\in\{0,1\}$，条件 $1\leqslant\sum_{i=1}^5|x_i|\leqslant 3$ 表示五个坐标中恰有 $1$，$2$ 或 $3$ 个非零。

\Needspace{5\baselineskip}
恰有 $k$ 个非零时，先选出这 $k$ 个位置有 $\mathrm C_5^k$ 种，每个非零坐标可取 $1$ 或 $-1$ 共 $2$ 种，故有 $\mathrm C_5^k2^k$ 个。于是
\nopagebreak[4]
\[\begin{aligned}
N&=\mathrm C_5^1\cdot2+\mathrm C_5^2\cdot2^2+\mathrm C_5^3\cdot2^3\\
&=5\times2+10\times4+10\times8\\
&=10+40+80=130.
\end{aligned}\]

选 D。

\gkanswer{答案：D。恰有 $1$，$2$，$3$ 个非零坐标的个数之和为 $\mathrm C_5^1\cdot2+\mathrm C_5^2\cdot4+\mathrm C_5^3\cdot8=10+40+80=130$。}

\end{answer}
\begin{answer}
\gksolhead{9}{}

\textbf{题目}：

不等式 $|x-1|+|x+2|\geqslant 5$ 的解集为\gkblank{}。

\Needspace{4\baselineskip}
\textbf{解析}

按绝对值内零点 $x=-2$，$x=1$ 分段。

\Needspace{5\baselineskip}
当 $x\leqslant -2$ 时，$|x-1|=1-x$，$|x+2|=-x-2$，不等式化为
\nopagebreak[4]
\[(1-x)+(-x-2)\geqslant 5\ \Longrightarrow\ -2x-1\geqslant 5\ \Longrightarrow\ x\leqslant -3,\]
与 $x\leqslant -2$ 取交得 $x\leqslant -3$。

当 $-2<x<1$ 时，$|x-1|=1-x$，$|x+2|=x+2$，左边恒为 $3$，$3\geqslant 5$ 不成立，无解。

\Needspace{5\baselineskip}
当 $x\geqslant 1$ 时，$|x-1|=x-1$，$|x+2|=x+2$，不等式化为
\nopagebreak[4]
\[2x+1\geqslant 5\ \Longrightarrow\ x\geqslant 2,\]
与 $x\geqslant 1$ 取交得 $x\geqslant 2$。

综上，解集为 $(-\infty,-3]\cup[2,+\infty)$。

\gkanswer{答案：$(-\infty,-3]\cup[2,+\infty)$。按 $x\leqslant-2$、$-2<x<1$、$x\geqslant1$ 分段去绝对值，分别得 $x\leqslant-3$、无解、$x\geqslant2$。}

\end{answer}
\begin{answer}
\gksolhead{10}{}

\textbf{题目}：

曲线 $y=\mathrm e^{-5x}+2$ 在点 $(0,3)$ 处的切线方程为\gkblank{}。

\Needspace{4\baselineskip}
\textbf{解析}

\Needspace{5\baselineskip}
$y=\mathrm e^{-5x}+2$ 的导数为
\nopagebreak[4]
\[y'=-5\mathrm e^{-5x}.\]
\Needspace{5\baselineskip}
在 $x=0$ 处 $y'=-5$，切点为 $(0,3)$，故切线方程为
\nopagebreak[4]
\[y-3=-5(x-0),\]
即 $y=-5x+3$，也可写成 $5x+y-3=0$。

\gkanswer{答案：$y=-5x+3$（即 $5x+y-3=0$）。$y'=-5\mathrm e^{-5x}$，在 $x=0$ 处斜率为 $-5$，过点 $(0,3)$。}

\end{answer}
\begin{answer}
\gksolhead{11}{}

\textbf{题目}：

从 $0,1,2,3,4,5,6,7,8,9$ 中任取七个不同的数，则这七个数的中位数是 $6$ 的概率为\gkblank{}。

\Needspace{4\baselineskip}
\textbf{解析}

从 $10$ 个数中任取 $7$ 个，共有 $\mathrm C_{10}^7=120$ 种取法。

把这 $7$ 个数从小到大排列，中位数是第 $4$ 个数。中位数为 $6$，等价于 $6$ 被取到，且所取数中恰有 $3$ 个小于 $6$、$3$ 个大于 $6$（$6$ 本身占第 $4$ 位）。

\Needspace{5\baselineskip}
小于 $6$ 的数有 $0,1,2,3,4,5$ 共 $6$ 个，取 $3$ 个有 $\mathrm C_6^3=20$ 种；大于 $6$ 的数有 $7,8,9$ 共 $3$ 个，取 $3$ 个只有 $\mathrm C_3^3=1$ 种。故有利取法为
\nopagebreak[4]
\[\mathrm C_6^3\mathrm C_3^3=20\times1=20.\]
\Needspace{5\baselineskip}
所求概率为
\nopagebreak[4]
\[P=\frac{20}{120}=\frac16.\]

\gkanswer{答案：$\dfrac16$。中位数为 $6$ 须恰取 $3$ 个小于 $6$ 的数和 $3$ 个大于 $6$ 的数，有利取法 $\mathrm C_6^3\mathrm C_3^3=20$，总取法 $\mathrm C_{10}^7=120$。}

\end{answer}
\begin{answer}
\gksolhead{12}{}

\textbf{题目}：

在 $\triangle ABC$ 中，角 $A$，$B$，$C$ 所对应的边分别为 $a$，$b$，$c$，已知 $b\cos C+c\cos B=2b$，则 $\dfrac ab=$\gkblank{}。

\Needspace{4\baselineskip}
\textbf{解析}

\Needspace{5\baselineskip}
由余弦定理，
\nopagebreak[4]
\[b\cos C=\frac{a^2+b^2-c^2}{2a},\qquad c\cos B=\frac{a^2+c^2-b^2}{2a},\]
\Needspace{5\baselineskip}
两式相加得
\nopagebreak[4]
\[b\cos C+c\cos B=\frac{2a^2}{2a}=a.\]
\Needspace{5\baselineskip}
（这正是投影公式：$b\cos C+c\cos B=a$。）代入已知条件 $b\cos C+c\cos B=2b$，得
\nopagebreak[4]
\[a=2b\ \Longrightarrow\ \frac ab=2.\]

\gkanswer{答案：$2$。由投影公式 $b\cos C+c\cos B=a$，得 $a=2b$，故 $\dfrac ab=2$。}

\end{answer}
\begin{answer}
\gksolhead{13}{}

\textbf{题目}：

若等比数列 $\{a_n\}$ 的各项均为正数，且 $a_{10}a_{11}+a_9a_{12}=2\mathrm e^5$，则 $\ln a_1+\ln a_2+\cdots+\ln a_{20}=$\gkblank{}。

\Needspace{4\baselineskip}
\textbf{解析}

\Needspace{5\baselineskip}
设公比为 $q$（$q>0$）。由等比数列性质，
\nopagebreak[4]
\[a_9a_{12}=a_{10}a_{11}=a_1^2q^{19},\]
\Needspace{5\baselineskip}
故已知条件化为
\nopagebreak[4]
\[2a_{10}a_{11}=2\mathrm e^5\ \Longrightarrow\ a_{10}a_{11}=\mathrm e^5.\]

\Needspace{5\baselineskip}
又对 $i=1,2,\ldots,10$，有
\nopagebreak[4]
\[a_i a_{21-i}=a_1^2q^{19}=a_{10}a_{11},\]
\Needspace{5\baselineskip}
即 $a_1a_{20}=a_2a_{19}=\cdots=a_{10}a_{11}=\mathrm e^5$。于是
\nopagebreak[4]
\[\prod_{i=1}^{20}a_i=(a_{10}a_{11})^{10}=\mathrm e^{50}.\]
\Needspace{5\baselineskip}
所以
\nopagebreak[4]
\[\ln a_1+\ln a_2+\cdots+\ln a_{20}
=\ln\prod_{i=1}^{20}a_i=\ln\mathrm e^{50}=50.\]

\gkanswer{答案：$50$。由 $a_9a_{12}=a_{10}a_{11}$ 得 $a_{10}a_{11}=\mathrm e^5$，而 $\prod_{i=1}^{20}a_i=(a_{10}a_{11})^{10}=\mathrm e^{50}$，取对数得 $50$。}

\end{answer}
\begin{answer}
\gksolhead{14}{}

\textbf{题目}：

在极坐标系中，曲线 $C_1$ 和 $C_2$ 的方程分别为 $\rho\sin^2\theta=\cos\theta$ 和 $\rho\sin\theta=1$，以极点为平面直角坐标系的原点，极轴为 $x$ 轴的正半轴，建立平面直角坐标系，则曲线 $C_1$ 和 $C_2$ 的交点的直角坐标为\gkblank{}。

\Needspace{4\baselineskip}
\textbf{解析}

由 $x=\rho\cos\theta$，$y=\rho\sin\theta$。

\Needspace{5\baselineskip}
对 $C_1$：$\rho\sin^2\theta=\cos\theta$，两边同乘 $\rho$ 得
\nopagebreak[4]
\[\rho^2\sin^2\theta=\rho\cos\theta\ \Longrightarrow\ y^2=x.\]
\Needspace{5\baselineskip}
对 $C_2$：$\rho\sin\theta=1$，即
\nopagebreak[4]
\[y=1.\]
把 $y=1$ 代入 $y^2=x$，得 $x=1$。故交点的直角坐标为 $(1,1)$。

（验证：点 $(1,1)$ 对应 $\rho=\sqrt2$，$\sin\theta=\cos\theta=\dfrac{\sqrt2}{2}$，满足 $\rho\sin^2\theta=\dfrac{\sqrt2}{2}=\cos\theta$ 且 $\rho\sin\theta=1$。）

\gkanswer{答案：$(1,1)$。$C_1$ 化为 $y^2=x$，$C_2$ 化为 $y=1$，联立得 $x=1$。}

\end{answer}
\begin{answer}
\gksolhead{15}{}

\textbf{题目}：

如图，在平行四边形 $ABCD$ 中，点 $E$ 在 $AB$ 上，且 $EB=2AE$，$AC$ 与 $DE$ 交于点 $F$，则 $\dfrac{\triangle CDF\ \text{\text{的面积}}}{\triangle AEF\ \text{\text{的面积}}}=$\gkblank{}。
\nopagebreak[4]
\begin{center}% ============================================================
% f2014_guangdong_science_q15_parallelogram.tex
% 2014 广东卷（理科）第 15 题：平行四边形 ABCD 中，E 在 AB 上，
% EB=2AE，DE 与 AC 交于点 F，求 △CDF 与 △AEF 的面积比。
% 来源 PDF：普通高考/2014/2014广东理.pdf
% 原页：第 1 页（A3 横排）第 2 栏，印面 877，第 15 题题干下方
% 说明：AE:EB=1:2，故 AE=AB/3；图中 A、E、B 依次在底边上，
% D 在 A 上方偏右，C 在 B 上方偏右，DE 近于竖直，AC 为对角线。
% ============================================================
\begin{tikzpicture}[>=Stealth, line width=0.7pt, font=\small,
  x=1cm, y=1cm, line cap=round, line join=round]

  \coordinate (A) at (0,0);
  \coordinate (E) at (1.3333,0);
  \coordinate (B) at (4,0);
  \coordinate (D) at (0.8,3);
  \coordinate (C) at (4.8,3);
  \coordinate (F) at (1.2,0.75);

  \draw (A) -- (B) -- (C) -- (D) -- cycle;
  \draw (D) -- (E);
  \draw (A) -- (C);

  \fill (A) circle (1.1pt) node[below left] {$A$};
  \fill (E) circle (1.1pt) node[below] {$E$};
  \fill (B) circle (1.1pt) node[below right] {$B$};
  \fill (D) circle (1.1pt) node[above left] {$D$};
  \fill (C) circle (1.1pt) node[above right] {$C$};
  \fill (F) circle (1.1pt) node[right] {$F$};
\end{tikzpicture}
\end{center}
\nopagebreak[4]
\Needspace{4\baselineskip}
\textbf{解析}

\Needspace{5\baselineskip}
在平行四边形 $ABCD$ 中，$AB\parallel DC$ 且 $AB=DC$，故 $AE\parallel DC$。于是
\nopagebreak[4]
\[\angle FAE=\angle FCD,\qquad \angle FEA=\angle FDC,\]
所以 $\triangle AEF\sim\triangle CDF$。

\Needspace{5\baselineskip}
由 $EB=2AE$ 得 $AB=AE+EB=3AE$，又 $DC=AB$，故相似比为
\nopagebreak[4]
\[\frac{AE}{DC}=\frac{AE}{AB}=\frac13.\]
\Needspace{5\baselineskip}
相似三角形的面积比等于相似比的平方，故
\nopagebreak[4]
\[\frac{S_{\triangle CDF}}{S_{\triangle AEF}}
=\left(\frac{DC}{AE}\right)^2=3^2=9.\]

\gkanswer{答案：$9$。$\triangle AEF\sim\triangle CDF$，相似比 $\dfrac{AE}{DC}=\dfrac13$，面积比为相似比的平方的倒数，即 $9$。}

\end{answer}
\begin{answer}
\gksolhead{16}{}

\textbf{题目}：

已知函数 $f(x)=A\sin\left(x+\dfrac\pi4\right)$，$x\in\mathbf R$，且 $f\left(\dfrac{5\pi}{12}\right)=\dfrac32$。

（1） 求 $A$ 的值；

（2） 若 $f(\theta)+f(-\theta)=\dfrac32$，$\theta\in\left(0,\dfrac\pi2\right)$，求 $f\left(\dfrac{3\pi}4-\theta\right)$。

\Needspace{4\baselineskip}
\textbf{解析}

\gksubpart{（1）}

\Needspace{5\baselineskip}
由 $f\left(\dfrac{5\pi}{12}\right)=\dfrac32$ 得
\nopagebreak[4]
\[A\sin\left(\frac{5\pi}{12}+\frac\pi4\right)=\frac32.\]
\Needspace{5\baselineskip}
而
\nopagebreak[4]
\[\frac{5\pi}{12}+\frac\pi4=\frac{5\pi}{12}+\frac{3\pi}{12}=\frac{8\pi}{12}=\frac{2\pi}{3},\]
\Needspace{5\baselineskip}
故
\nopagebreak[4]
\[A\sin\frac{2\pi}{3}=A\cdot\frac{\sqrt3}{2}=\frac32\ \Longrightarrow\ A=\sqrt3.\]

\gksubpart{（2）}

\Needspace{5\baselineskip}
由（1），$f(x)=\sqrt3\sin\left(x+\dfrac\pi4\right)$。于是
\nopagebreak[4]
\[\begin{aligned}
f(\theta)+f(-\theta)
&=\sqrt3\left[\sin\left(\theta+\frac\pi4\right)+\sin\left(-\theta+\frac\pi4\right)\right]\\
&=\sqrt3\cdot2\sin\frac\pi4\cos\theta
=\sqrt3\cdot\sqrt2\cos\theta=\sqrt6\cos\theta.
\end{aligned}\]
\Needspace{5\baselineskip}
由 $f(\theta)+f(-\theta)=\dfrac32$ 得
\nopagebreak[4]
\[\cos\theta=\frac{3}{2\sqrt6}=\frac{\sqrt6}{4}.\]
\Needspace{5\baselineskip}
因 $\theta\in\left(0,\dfrac\pi2\right)$，$\sin\theta>0$，故
\nopagebreak[4]
\[\sin\theta=\sqrt{1-\cos^2\theta}
=\sqrt{1-\frac{6}{16}}=\sqrt{\frac{10}{16}}=\frac{\sqrt{10}}{4}.\]
\Needspace{5\baselineskip}
所求为
\nopagebreak[4]
\[\begin{aligned}
f\left(\frac{3\pi}{4}-\theta\right)
&=\sqrt3\sin\left(\frac{3\pi}{4}-\theta+\frac\pi4\right)\\
&=\sqrt3\sin(\pi-\theta)=\sqrt3\sin\theta
=\sqrt3\cdot\frac{\sqrt{10}}{4}=\frac{\sqrt{30}}{4}.
\end{aligned}\]

\gkanswer{答案（1）：$\sqrt3$。$\sin\left(\dfrac{5\pi}{12}+\dfrac\pi4\right)=\sin\dfrac{2\pi}{3}=\dfrac{\sqrt3}{2}$，故 $A\cdot\dfrac{\sqrt3}{2}=\dfrac32$。}

\gkanswer{答案（2）：$\dfrac{\sqrt{30}}4$。$f(\theta)+f(-\theta)=\sqrt6\cos\theta=\dfrac32$ 得 $\cos\theta=\dfrac{\sqrt6}{4}$，$\sin\theta=\dfrac{\sqrt{10}}4$，所求为 $\sqrt3\sin\theta$。}

\end{answer}
\begin{answer}
\gksolhead{17}{}

\textbf{题目}：

随机观测生产某种零件的某工厂 $25$ 名工人的日加工零件数（单位：件），获得数据如下：$30$，$42$，$41$，$36$，$44$，$40$，$37$，$37$，$25$，$45$，$29$，$43$，$31$，$36$，$49$，$34$，$33$，$43$，$38$，$42$，$32$，$34$，$46$，$39$，$36$。根据上述数据得到样本的频率分布表如下：

\begin{center}
\begin{tabular}{lll}
\toprule
分组 & 频数 & 频率 \\
\midrule
$[25,30]$ & $3$ & $0.12$ \\
$(30,35]$ & $5$ & $0.20$ \\
$(35,40]$ & $8$ & $0.32$ \\
$(40,45]$ & $n_1$ & $f_1$ \\
$(45,50]$ & $n_2$ & $f_2$ \\
\bottomrule
\end{tabular}
\end{center}

（1） 确定样本频率分布表中 $n_1$，$n_2$，$f_1$ 和 $f_2$ 的值；

（2） 根据上述频率分布表，画出样本频率分布直方图；

（3） 根据样本频率分布直方图，求在该厂任取 $4$ 人，至少有 $1$ 人的日加工零件数落在区间 $(30,35]$ 的概率。

\Needspace{4\baselineskip}
\textbf{解析}

\gksubpart{（1）}

把 $25$ 个数据按区间计数：

$[25,30]$ 内有 $25,29,30$，共 $3$ 个；

$(30,35]$ 内有 $31,32,33,34,34$，共 $5$ 个；

$(35,40]$ 内有 $36,36,36,37,37,38,39,40$，共 $8$ 个；

$(40,45]$ 内有 $41,42,42,43,43,44$，共 $6$ 个；

$(45,50]$ 内有 $45,46,49$，共 $3$ 个。

\Needspace{5\baselineskip}
合计 $3+5+8+6+3=25$，与总人数相符。故
\nopagebreak[4]
\[n_1=6,\qquad n_2=3,\]
\nopagebreak[4]
\[f_1=\frac{6}{25}=0.24,\qquad f_2=\frac{3}{25}=0.12.\]

\gksubpart{（2）}

\Needspace{5\baselineskip}
频率分布直方图的纵轴是 $\dfrac{\text{\text{频率}}}{\text{\text{组距}}}$，组距为 $5$。各矩形的高为
\nopagebreak[4]
\[\begin{aligned}
[25,30]&:\ \frac{0.12}{5}=0.024,\\
(30,35]&:\ \frac{0.20}{5}=0.040,\\
(35,40]&:\ \frac{0.32}{5}=0.064,\\
(40,45]&:\ \frac{0.24}{5}=0.048,\\
(45,50]&:\ \frac{0.12}{5}=0.024.
\end{aligned}\]
以分组为横轴、上述高为纵轴，画出五个等宽矩形，即得样本频率分布直方图。

\gksubpart{（3）}

\Needspace{5\baselineskip}
由频率分布表，日加工零件数落在 $(30,35]$ 的频率为 $0.20$，故任取 $1$ 人落在该区间的概率为
\nopagebreak[4]
\[p=0.2=\frac15.\]
\Needspace{5\baselineskip}
任取 $4$ 人，设落在 $(30,35]$ 的人数为 $X$，则 $X\sim B\left(4,\dfrac15\right)$。“至少有 $1$ 人”的对立事件是“$4$ 人都不落在该区间”，故
\nopagebreak[4]
\[\begin{aligned}
P(X\geqslant 1)&=1-P(X=0)\\
&=1-\left(1-\frac15\right)^4
=1-\left(\frac45\right)^4\\
&=1-\frac{256}{625}=\frac{369}{625}\approx 0.5904.
\end{aligned}\]

\gkanswer{答案（1）：$n_1=6$，$n_2=3$，$f_1=0.24$，$f_2=0.12$。$(40,45]$ 内有 $6$ 个数据，$(45,50]$ 内有 $3$ 个，频率分别为 $\dfrac6{25}$ 与 $\dfrac3{25}$。}

\gkanswer{答案（2）：纵轴取 $\dfrac{\text{\text{频率}}}{\text{\text{组距}}}$，五个矩形的高依次为 $0.024$，$0.040$，$0.064$，$0.048$，$0.024$。组距为 $5$，各高等于对应频率除以 $5$。}

\gkanswer{答案（3）：$\dfrac{369}{625}$（约 $0.5904$）。落在 $(30,35]$ 的概率为 $\dfrac15$，任取 $4$ 人至少 $1$ 人落在该区间的概率为 $1-\left(\dfrac45\right)^4$。}

\end{answer}
\begin{answer}
\gksolhead{18}{}

\textbf{题目}：

如图，四边形 $ABCD$ 为正方形，$PD\perp$ 平面 $ABCD$，$\angle DPC=30^\circ$，$AF\perp PC$ 于点 $F$，$FE\parallel CD$，交 $PD$ 于点 $E$。
\nopagebreak[4]
\begin{center}% ============================================================
% f2014_guangdong_science_q18_pyramid.tex
% 2014 广东卷（理科）第 18 题：四边形 ABCD 为正方形，
% PD⊥平面 ABCD，∠DPC=30°，AF⊥PC 于 F，FE∥CD 交 PD 于 E。
% 来源 PDF：普通高考/2014/2014广东理.pdf
% 原页：第 1 页（A3 横排）第 3 栏下部，印面 877，第 18 题题干下方
% 说明：底面 ABCD 画成平行四边形以示意正方形；PD 为自 D 向下的
% 棱，E 在 PD 上，F 在 PC 上，EF∥CD；被底面遮挡的 AD、DC、DF
% 用虚线，其余可见棱用实线。
% ============================================================
\begin{tikzpicture}[>=Stealth, line width=0.7pt, font=\small,
  x=1cm, y=1cm, line cap=round, line join=round]

  \coordinate (A) at (0,3.4);
  \coordinate (B) at (4.2,3.4);
  \coordinate (C) at (5.2,0.9);
  \coordinate (D) at (1.0,0.9);
  \coordinate (P) at (-0.4,-1.8);
  \coordinate (E) at (0.65,0.225);
  \coordinate (F) at (3.8,0.225);

  % 被遮挡的棱
  \draw[dashed] (A) -- (D);
  \draw[dashed] (D) -- (C);
  \draw[dashed] (D) -- (F);

  % 可见的棱
  \draw (A) -- (B) -- (C);
  \draw (P) -- (E) -- (D);
  \draw (P) -- (C);
  \draw (E) -- (F);
  \draw (A) -- (C);
  \draw (A) -- (F);

  \fill (A) circle (1.1pt) node[above left] {$A$};
  \fill (B) circle (1.1pt) node[above right] {$B$};
  \fill (C) circle (1.1pt) node[right] {$C$};
  \fill (D) circle (1.1pt) node[above right, inner sep=2pt] {$D$};
  \fill (P) circle (1.1pt) node[below left] {$P$};
  \fill (E) circle (1.1pt) node[left] {$E$};
  \fill (F) circle (1.1pt) node[below right] {$F$};
\end{tikzpicture}
\end{center}
\nopagebreak[4]
（1） 证明：$CF\perp$ 平面 $ADF$；

（2） 求二面角 $D-AF-E$ 的余弦值。

\Needspace{4\baselineskip}
\textbf{解析}

\gksubpart{（1）}

\Needspace{5\baselineskip}
因为 $PD\perp$ 平面 $ABCD$，$AD\subset$ 平面 $ABCD$，所以
\nopagebreak[4]
\[PD\perp AD.\]
\Needspace{5\baselineskip}
又 $ABCD$ 为正方形，$AD\perp DC$，且 $PD\cap DC=D$，故
\nopagebreak[4]
\[AD\perp\text{\text{平面} }PDC.\]
\Needspace{5\baselineskip}
因为 $CF\subset$ 平面 $PDC$，所以
\nopagebreak[4]
\[AD\perp CF.\]
\Needspace{5\baselineskip}
又已知 $AF\perp PC$，即 $AF\perp CF$，且 $AD\cap AF=A$，$AD$，$AF\subset$ 平面 $ADF$，所以
\nopagebreak[4]
\[CF\perp\text{\text{平面} }ADF.\]

\gksubpart{（2）}

\Needspace{5\baselineskip}
设正方形边长为 $1$。因为 $PD\perp$ 平面 $ABCD$，$DC\subset$ 平面 $ABCD$，所以 $PD\perp DC$，$\triangle PDC$ 在 $D$ 处为直角。由 $\angle DPC=30^\circ$ 得
\nopagebreak[4]
\[\tan 30^\circ=\frac{DC}{PD}=\frac1{PD}
\ \Longrightarrow\ PD=\sqrt3.\]

\Needspace{5\baselineskip}
建立空间直角坐标系：以 $D$ 为原点，$\overrightarrow{DA}$ 为 $x$ 轴正方向，$\overrightarrow{DC}$ 为 $y$ 轴正方向，$\overrightarrow{DP}$ 为 $z$ 轴正方向。则
\nopagebreak[4]
\[D(0,0,0),\quad A(1,0,0),\quad C(0,1,0),\quad P(0,0,\sqrt3).\]

\Needspace{5\baselineskip}
设 $F=P+t(\overrightarrow{PC})$，$0<t<1$，即
\nopagebreak[4]
\[F=(0,t,\sqrt3(1-t)).\]
\Needspace{5\baselineskip}
则 $\overrightarrow{AF}=(-1,t,\sqrt3(1-t))$，$\overrightarrow{PC}=(0,1,-\sqrt3)$。由 $AF\perp PC$ 得
\nopagebreak[4]
\[\overrightarrow{AF}\cdot\overrightarrow{PC}=t-3(1-t)=4t-3=0
\ \Longrightarrow\ t=\frac34.\]
\Needspace{5\baselineskip}
故
\nopagebreak[4]
\[F=\left(0,\frac34,\frac{\sqrt3}{4}\right).\]

\Needspace{5\baselineskip}
因为 $FE\parallel CD$，而 $\overrightarrow{CD}=(0,-1,0)$，从 $F$ 沿 $(0,-1,0)$ 方向移动到 $PD$（即 $x=y=0$）需移动 $\dfrac34$，得
\nopagebreak[4]
\[E=\left(0,0,\frac{\sqrt3}{4}\right).\]

\Needspace{5\baselineskip}
平面 $ADF$ 的一个法向量可取 $\overrightarrow{CF}$：
\nopagebreak[4]
\[\overrightarrow{CF}=\left(0,-\frac14,\frac{\sqrt3}{4}\right),\]
与其同向的 $\boldsymbol n_1=(0,-1,\sqrt3)$ 也是法向量。

\Needspace{5\baselineskip}
平面 $AEF$ 的法向量为
\nopagebreak[4]
\[\begin{aligned}
\boldsymbol n_2
&=\overrightarrow{AE}\times\overrightarrow{AF}\\
&=\left(-1,0,\frac{\sqrt3}{4}\right)\times\left(-1,\frac34,\frac{\sqrt3}{4}\right)\\
&=\left(-\frac{3\sqrt3}{16},\,0,\,-\frac34\right),
\end{aligned}\]
与其同向的 $\boldsymbol n_2'=(\sqrt3,0,4)$ 也是法向量（各分量乘以 $-\dfrac{16}{3}$）。

\Needspace{5\baselineskip}
于是
\nopagebreak[4]
\[\begin{aligned}
|\cos\langle\boldsymbol n_1,\boldsymbol n_2'\rangle|
&=\frac{|\boldsymbol n_1\cdot\boldsymbol n_2'|}
{|\boldsymbol n_1|\,|\boldsymbol n_2'|}\\
&=\frac{|0\cdot\sqrt3+(-1)\cdot0+\sqrt3\cdot4|}
{\sqrt{0+1+3}\cdot\sqrt{3+0+16}}\\
&=\frac{4\sqrt3}{2\sqrt{19}}
=\frac{2\sqrt3}{\sqrt{19}}
=\frac{2\sqrt{57}}{19}.
\end{aligned}\]
二面角 $D-AF-E$ 为锐角（结合图形，两半平面所张的角小于 $90^\circ$），故其余弦值为 $\dfrac{2\sqrt{57}}{19}$。

\gkanswer{答案（1）：见解析版}

\gkanswer{答案（2）：$\dfrac{2\sqrt{57}}{19}$。设边长 $1$，$PD=\sqrt3$，取 $D$ 为原点建系得 $F\left(0,\dfrac34,\dfrac{\sqrt3}{4}\right)$，$E\left(0,0,\dfrac{\sqrt3}{4}\right)$，两平面法向量为 $(0,-1,\sqrt3)$ 与 $(\sqrt3,0,4)$。}

\end{answer}
\begin{answer}
\gksolhead{19}{}

\textbf{题目}：

设数列 $\{a_n\}$ 的前 $n$ 项和为 $S_n$，满足 $S_n=2na_{n+1}-3n^2-4n$，$n\in\mathbf N^*$，且 $S_3=15$。

（1） 求 $a_1$，$a_2$，$a_3$ 的值；

（2） 求数列 $\{a_n\}$ 的通项公式。

\Needspace{4\baselineskip}
\textbf{解析}

\gksubpart{（1）}

在 $S_n=2na_{n+1}-3n^2-4n$ 中依次取 $n=1,2,3$。

\Needspace{5\baselineskip}
$n=1$：$S_1=a_1=2a_2-3-4$，即
\nopagebreak[4]
\[a_1=2a_2-7. \tag{1}\]

\Needspace{5\baselineskip}
$n=2$：$S_2=a_1+a_2=4a_3-12-8$，即
\nopagebreak[4]
\[a_1+a_2=4a_3-20. \tag{2}\]

$n=3$：$S_3=a_1+a_2+a_3=6a_4-27-12$，结合 $S_3=15$ 只用于验证，暂不参与求 $a_1,a_2,a_3$。

\Needspace{5\baselineskip}
由 $S_3=15$ 得
\nopagebreak[4]
\[a_3=15-(a_1+a_2).\]
\Needspace{5\baselineskip}
代入（2）：
\nopagebreak[4]
\[a_3=15-(4a_3-20)=35-4a_3
\ \Longrightarrow\ 5a_3=35
\ \Longrightarrow\ a_3=7.\]
\Needspace{5\baselineskip}
于是 $a_1+a_2=4\times7-20=8$。与（1）联立：
\nopagebreak[4]
\[(2a_2-7)+a_2=8
\ \Longrightarrow\ 3a_2=15
\ \Longrightarrow\ a_2=5,\]
\nopagebreak[4]
\[a_1=2\times5-7=3.\]
\Needspace{5\baselineskip}
故
\nopagebreak[4]
\[a_1=3,\qquad a_2=5,\qquad a_3=7.\]

\gksubpart{（2）}

\Needspace{5\baselineskip}
由 $a_{n+1}=S_{n+1}-S_n$，把 $S_n=2na_{n+1}-3n^2-4n$ 改写为
\nopagebreak[4]
\[S_n=2n(S_{n+1}-S_n)-3n^2-4n,\]
\Needspace{5\baselineskip}
即
\nopagebreak[4]
\[(2n+1)S_n=2nS_{n+1}-3n^2-4n,\]
\Needspace{5\baselineskip}
所以
\nopagebreak[4]
\[S_{n+1}=\frac{(2n+1)S_n+3n^2+4n}{2n}. \tag{3}\]

\Needspace{5\baselineskip}
由（1）$S_1=3$。用（3）算前几项：
\nopagebreak[4]
\[S_2=\frac{3\cdot3+3+4}{2}=8,\qquad
S_3=\frac{5\cdot8+12+8}{4}=15,\]
与题设 $S_3=15$ 一致。

\Needspace{5\baselineskip}
猜测 $S_n=n^2+2n$。下面用数学归纳法证明。$n=1$ 时 $S_1=3=1^2+2$，成立。设 $S_n=n^2+2n$，则由（3）
\nopagebreak[4]
\[\begin{aligned}
S_{n+1}
&=\frac{(2n+1)(n^2+2n)+3n^2+4n}{2n}\\
&=\frac{(2n+1)n(n+2)+n(3n+4)}{2n}\\
&=\frac{(2n+1)(n+2)+3n+4}{2}\\
&=\frac{2n^2+5n+2+3n+4}{2}\\
&=\frac{2n^2+8n+6}{2}
=n^2+4n+3\\
&=(n+1)^2+2(n+1),
\end{aligned}\]
\Needspace{5\baselineskip}
归纳成立。故
\nopagebreak[4]
\[S_n=n^2+2n\quad(n\in\mathbf N^*).\]

\Needspace{5\baselineskip}
当 $n\geqslant2$ 时，
\nopagebreak[4]
\[a_n=S_n-S_{n-1}=(n^2+2n)-[(n-1)^2+2(n-1)]=2n+1.\]
\Needspace{5\baselineskip}
当 $n=1$ 时，$a_1=3=2\times1+1$，也符合。所以
\nopagebreak[4]
\[a_n=2n+1\quad(n\in\mathbf N^*).\]

（验证：$S_n=\sum_{k=1}^n(2k+1)=n(n+2)$，且 $2na_{n+1}-3n^2-4n=2n(2n+3)-3n^2-4n=n^2+2n=S_n$，满足题设。）

\gkanswer{答案（1）：$a_1=3$，$a_2=5$，$a_3=7$。由 $n=1,2$ 及 $S_3=15$ 联立解得。}

\gkanswer{答案（2）：$a_n=2n+1$。由递推式得 $S_{n+1}=\dfrac{(2n+1)S_n+3n^2+4n}{2n}$，归纳得 $S_n=n^2+2n$，相减得 $a_n=2n+1$。}

\end{answer}
\begin{answer}
\gksolhead{20}{}

\textbf{题目}：

已知椭圆 $C:\dfrac{x^2}{a^2}+\dfrac{y^2}{b^2}=1$（$a>b>0$）的一个焦点为 $(\sqrt5,0)$，离心率为 $\dfrac{\sqrt5}{3}$。

（1） 求椭圆 $C$ 的标准方程；

（2） 若动点 $P(x_0,y_0)$ 为椭圆 $C$ 外一点，且点 $P$ 到椭圆 $C$ 的两条切线相互垂直，求点 $P$ 的轨迹方程。

\Needspace{4\baselineskip}
\textbf{解析}

\gksubpart{（1）}

\Needspace{5\baselineskip}
设焦距为 $2c$。由焦点 $(\sqrt5,0)$ 得 $c=\sqrt5$。由离心率
\nopagebreak[4]
\[e=\frac ca=\frac{\sqrt5}{3}\]
\Needspace{5\baselineskip}
得
\nopagebreak[4]
\[a=\frac{c}{e}=\frac{\sqrt5}{\sqrt5/3}=3.\]
\Needspace{5\baselineskip}
于是
\nopagebreak[4]
\[b^2=a^2-c^2=9-5=4.\]
\Needspace{5\baselineskip}
故椭圆 $C$ 的标准方程为
\nopagebreak[4]
\[\frac{x^2}{9}+\frac{y^2}{4}=1.\]

\gksubpart{（2）}

\Needspace{5\baselineskip}
先设过 $P(x_0,y_0)$ 的切线斜率存在，记为 $k$，切线为
\nopagebreak[4]
\[y=kx+m,\qquad m=y_0-kx_0.\]
\Needspace{5\baselineskip}
代入椭圆方程 $\dfrac{x^2}{9}+\dfrac{y^2}{4}=1$，整理得
\nopagebreak[4]
\[(4+9k^2)x^2+18km\,x+9m^2-36=0.\]
\Needspace{5\baselineskip}
切线与椭圆只有一个公共点，判别式为零：
\nopagebreak[4]
\[(18km)^2-4(4+9k^2)(9m^2-36)=0.\]
\Needspace{5\baselineskip}
化简：
\nopagebreak[4]
\[\begin{aligned}
324k^2m^2-36(4+9k^2)(m^2-4)&=0\\
9k^2m^2-(4+9k^2)(m^2-4)&=0\\
9k^2m^2-4m^2+16-9k^2m^2+36k^2&=0\\
-4m^2+16+36k^2&=0,
\end{aligned}\]
\Needspace{5\baselineskip}
即
\nopagebreak[4]
\[m^2=4+9k^2.\]
\Needspace{5\baselineskip}
代入 $m=y_0-kx_0$：
\nopagebreak[4]
\[(y_0-kx_0)^2=4+9k^2,\]
\Needspace{5\baselineskip}
整理成关于 $k$ 的二次方程
\nopagebreak[4]
\[(x_0^2-9)k^2-2x_0y_0k+(y_0^2-4)=0. \tag{4}\]
\Needspace{5\baselineskip}
两条切线的斜率 $k_1$，$k_2$ 是（4）的两根。由两切线相互垂直得
\nopagebreak[4]
\[k_1k_2=-1.\]
\Needspace{5\baselineskip}
由韦达定理
\nopagebreak[4]
\[k_1k_2=\frac{y_0^2-4}{x_0^2-9},\]
\Needspace{5\baselineskip}
故
\nopagebreak[4]
\[\frac{y_0^2-4}{x_0^2-9}=-1
\ \Longrightarrow\ y_0^2-4=9-x_0^2
\ \Longrightarrow\ x_0^2+y_0^2=13. \tag{5}\]

再补上一条切线竖直的情形。若 $x_0=\pm3$，则 $x=x_0$ 是椭圆的一条切线；要与它垂直，另一条切线必须是水平线 $y=y_0$，这要求 $y_0=\pm2$。此时 $P=(\pm3,\pm2)$，仍有 $x_0^2+y_0^2=9+4=13$，即（5）依旧成立。

\Needspace{5\baselineskip}
反过来，若 $P(x_0,y_0)$ 满足 $x_0^2+y_0^2=13$，则
\nopagebreak[4]
\[\frac{x_0^2}{9}+\frac{y_0^2}{4}
=\frac{x_0^2}{9}+\frac{13-x_0^2}{4}
=\frac{13}{4}+\frac{5x_0^2}{36}\geqslant\frac{13}{4}>1,\]
\Needspace{5\baselineskip}
故 $P$ 确在椭圆外；且由（4）知两条切线的斜率之积为 $-1$（竖直切线情形同样满足）。因此点 $P$ 的轨迹方程为
\nopagebreak[4]
\[x^2+y^2=13.\]

\gkanswer{答案（1）：$\dfrac{x^2}{9}+\dfrac{y^2}{4}=1$。$c=\sqrt5$，由 $\dfrac ca=\dfrac{\sqrt5}{3}$ 得 $a=3$，$b^2=9-5=4$。}

\gkanswer{答案（2）：$x^2+y^2=13$。设切线 $y=kx+m$，相切条件化为 $(x_0^2-9)k^2-2x_0y_0k+(y_0^2-4)=0$，由 $k_1k_2=-1$ 得 $x_0^2+y_0^2=13$（竖直切线情形同样满足）。}

\end{answer}
\begin{answer}
\gksolhead{21}{}

\textbf{题目}：

设函数 $f(x)=\dfrac1{\sqrt{(x^2+2x+k)^2+2(x^2+2x+k)-3}}$，其中 $k<-2$。

（1） 求函数 $f(x)$ 的定义域 $D$（用区间表示）；

（2） 讨论 $f(x)$ 在区间 $D$ 上的单调性；

（3） 若 $k<-6$，求 $D$ 上满足条件 $f(x)>f(1)$ 的 $x$ 的集合（用区间表示）。

\Needspace{4\baselineskip}
\textbf{解析}

\Needspace{5\baselineskip}
记
\nopagebreak[4]
\[t=x^2+2x+k=(x+1)^2+k-1.\]
\Needspace{5\baselineskip}
则
\nopagebreak[4]
\[f(x)=\frac1{\sqrt{t^2+2t-3}}=\frac1{\sqrt{(t+3)(t-1)}}.\]

\gksubpart{（1）}

\Needspace{5\baselineskip}
定义域要求根号内为正：
\nopagebreak[4]
\[t^2+2t-3>0\ \Longleftrightarrow\ (t+3)(t-1)>0
\ \Longleftrightarrow\ t<-3\ \text{\text{或}}\ t>1.\]

由 $t=(x+1)^2+k-1$：

\Needspace{5\baselineskip}
$t>1$ 等价于 $(x+1)^2>2-k$，即
\nopagebreak[4]
\[x<-1-\sqrt{2-k}\quad\text{\text{或}}\quad x>-1+\sqrt{2-k}.\]

\Needspace{5\baselineskip}
$t<-3$ 等价于 $(x+1)^2<-2-k$。因 $k<-2$，有 $-2-k>0$，解得
\nopagebreak[4]
\[-1-\sqrt{-2-k}<x<-1+\sqrt{-2-k}.\]

\Needspace{5\baselineskip}
故
\nopagebreak[4]
\[D=(-\infty,-1-\sqrt{2-k})\cup(-1-\sqrt{-2-k},-1+\sqrt{-2-k})\cup(-1+\sqrt{2-k},+\infty).\]

\gksubpart{（2）}

\Needspace{5\baselineskip}
由 $f(x)=(t^2+2t-3)^{-1/2}$ 及 $\dfrac{\mathrm dt}{\mathrm dx}=2x+2$，得
\nopagebreak[4]
\[\begin{aligned}
f'(x)
&=-\frac12(t^2+2t-3)^{-3/2}\cdot(2t+2)(2x+2)\\
&=-\frac{2(t+1)(x+1)}{(t^2+2t-3)^{3/2}}.
\end{aligned}\]
分母恒正，故 $f'(x)$ 的符号由 $-(t+1)(x+1)$ 决定。

在 $t>1$ 的两段上，$t+1>0$：

当 $x<-1-\sqrt{2-k}$ 时 $x+1<0$，$(t+1)(x+1)<0$，故 $f'(x)>0$；

当 $x>-1+\sqrt{2-k}$ 时 $x+1>0$，$(t+1)(x+1)>0$，故 $f'(x)<0$。

在 $t<-3$ 的一段 $(-1-\sqrt{-2-k},-1+\sqrt{-2-k})$ 上，$t+1<0$：

当 $-1-\sqrt{-2-k}<x<-1$ 时 $x+1<0$，$(t+1)(x+1)>0$，故 $f'(x)<0$；

当 $-1<x<-1+\sqrt{-2-k}$ 时 $x+1>0$，$(t+1)(x+1)<0$，故 $f'(x)>0$。

$x=-1$ 时 $f'(x)=0$，且 $-1\in D$。

\Needspace{5\baselineskip}
综上，$f(x)$ 在 $D$ 上的单调性为：
\nopagebreak[4]
\[\text{\text{增区间：}}\ (-\infty,-1-\sqrt{2-k}),\quad(-1,-1+\sqrt{-2-k});\]
\nopagebreak[4]
\[\text{\text{减区间：}}\ (-1-\sqrt{-2-k},-1),\quad(-1+\sqrt{2-k},+\infty).\]

\gksubpart{（3）}

\Needspace{5\baselineskip}
设 $k<-6$。此时
\nopagebreak[4]
\[t_0=1^2+2\times1+k=k+3<-3,\]
\Needspace{5\baselineskip}
故 $1\in D$，且
\nopagebreak[4]
\[f(1)=\frac1{\sqrt{(k+3)^2+2(k+3)-3}}.\]
\Needspace{5\baselineskip}
记 $g(t)=t^2+2t-3=(t+1)^2-4$。因 $f(x)>0$，$f(1)>0$，故
\nopagebreak[4]
\[f(x)>f(1)\ \Longleftrightarrow\ g(t)<g(t_0).\]
\Needspace{5\baselineskip}
而
\nopagebreak[4]
\[g(t)<g(t_0)\ \Longleftrightarrow\ (t+1)^2<(t_0+1)^2.\]
\Needspace{5\baselineskip}
由 $t_0+1=k+4<-2$，得 $|t_0+1|=-(k+4)=-k-4$。又
\nopagebreak[4]
\[t+1=(x+1)^2+k,\]
\Needspace{5\baselineskip}
故条件化为
\nopagebreak[4]
\[|(x+1)^2+k|<-k-4. \tag{6}\]

\Needspace{5\baselineskip}
令 $u=(x+1)^2\geqslant0$，则（6）即
\nopagebreak[4]
\[-(-k-4)<u+k<-k-4
\ \Longleftrightarrow\ 4<u<-2k-4. \tag{7}\]
\Needspace{5\baselineskip}
同时 $x\in D$ 要求 $g(t)>0$，即 $|t+1|>2$，亦即
\nopagebreak[4]
\[|u+k|>2. \tag{8}\]

\Needspace{5\baselineskip}
由 $k<-6$ 得
\nopagebreak[4]
\[4<-k-2<-k+2<-2k-4.\]
\Needspace{5\baselineskip}
在（7）的 $u$ 范围内，条件（8）等价于
\nopagebreak[4]
\[u<-k-2\quad\text{\text{或}}\quad u>-k+2,\]
\Needspace{5\baselineskip}
与（7）取交得
\nopagebreak[4]
\[u\in(4,-k-2)\cup(-k+2,-2k-4).\]

代回 $u=(x+1)^2$：

\Needspace{5\baselineskip}
由 $4<(x+1)^2<-k-2$ 得
\nopagebreak[4]
\[-1-\sqrt{-k-2}<x<-3
\quad\text{\text{或}}\quad
1<x<-1+\sqrt{-k-2};\]

\Needspace{5\baselineskip}
由 $-k+2<(x+1)^2<-2k-4$ 得
\nopagebreak[4]
\[-1-\sqrt{-2k-4}<x<-1-\sqrt{-k+2}
\quad\text{\text{或}}\quad
-1+\sqrt{-k+2}<x<-1+\sqrt{-2k-4}.\]

\Needspace{5\baselineskip}
故所求 $x$ 的集合为
\nopagebreak[4]
\[\begin{aligned}
&\left(-1-\sqrt{-2k-4},\,-1-\sqrt{-k+2}\right)\\
&\cup\left(-1-\sqrt{-k-2},\,-3\right)\\
&\cup\left(1,\,-1+\sqrt{-k-2}\right)\\
&\cup\left(-1+\sqrt{-k+2},\,-1+\sqrt{-2k-4}\right).
\end{aligned}\]

\gkanswer{答案（1）：$D=(-\infty,-1-\sqrt{2-k})\cup(-1-\sqrt{-2-k},-1+\sqrt{-2-k})\cup(-1+\sqrt{2-k},+\infty)$。令 $t=x^2+2x+k$，由 $t^2+2t-3>0$ 得 $t<-3$ 或 $t>1$，再解出 $x$。}

\gkanswer{答案（2）：增区间为 $(-\infty,-1-\sqrt{2-k})$ 与 $(-1,-1+\sqrt{-2-k})$；减区间为 $(-1-\sqrt{-2-k},-1)$ 与 $(-1+\sqrt{2-k},+\infty)$。由 $f'(x)=-\dfrac{2(t+1)(x+1)}{(t^2+2t-3)^{3/2}}$ 定号。}

\gkanswer{答案（3）：$\left(-1-\sqrt{-2k-4},-1-\sqrt{-k+2}\right)\cup\left(-1-\sqrt{-k-2},-3\right)\cup\left(1,-1+\sqrt{-k-2}\right)\cup\left(-1+\sqrt{-k+2},-1+\sqrt{-2k-4}\right)$。由 $f(x)>f(1)$ 化为 $|(x+1)^2+k|<-k-4$，再与定义域条件 $|(x+1)^2+k|>2$ 取交。}

\end{answer}